\documentclass{article}

% Margins and spacing
\usepackage[a4paper, margin=.7in]{geometry}
% Math packages
\usepackage{amsmath}
\usepackage{amssymb}
% Placement of figures, tables, algorithms, etc.
\usepackage{float}
% Algorithms
\usepackage{algorithm}
\usepackage{algpseudocode}

% No section numbering
\setcounter{secnumdepth}{0}
% No paragraph indentation
\setlength{\parindent}{0pt}

% Define custom commands
%% Operators
\DeclareMathOperator*{\argmin}{arg\,min}
\DeclareMathOperator*{\argmax}{arg\,max}
%% Manually indented state in algorithmic environment
\newcommand{\IndentState}{\State \hspace{2em}}

\title{CSCI-570 Homework 2 Solutions}
\author{
  Naman Chhibbar \\
  USC ID: 7709-8548-02
}
\date{}

% Document body
\begin{document}
\maketitle

\section{Problem 1}
  Online insertion of $n$ elements into a binary heap using an unbounded array can be broken into online insertion of $n$ elements into a binary heap and into an unbounded array of starting size 1 independently.
  \begin{itemize}
    \item \textbf{Insertion into a binary heap:} We know from the lectures that inserting an element into a binary heap takes $O(\log n)$ time, where $n$ is the number of elements in the heap.
    Hence, inserting $n$ elements will take:
    \begin{align*}
      T(n) &= \sum_{i=1}^{n} O(\log i) \\
      &\le \sum_{i=1}^{n} O(\log n) \\
      &= O(n \log n)
    \end{align*}

    \item \textbf{Insertion into an unbounded array:} We start with an array of size 1.
    Addition of $n$ elements into the array (without considering the time taken to copy elements when the array is full) takes $O(n)$ time.
    Now, for simplicity, assume that $n = 2^k$ for some $k \in \mathbb{N}$ (we can do this since we are analyzing the time complexity asymptotically).
    We can see that we will need to resize the array $k$ times.
    Hence, the total time taken will be:
    \begin{align*}
      T(n) &= O(n) + \sum_{i=0}^{k-1} 2^i \\
      &= O(n) + (2^k - 1) \\
      &= O(n) + (n - 1) \\
      &= O(n)
    \end{align*}
  \end{itemize}
  Therefore, the total time taken to insert $n$ elements into a binary heap using an unbounded array is $O(n \log n) + O(n) = O(n \log n)$.

\section{Problem 2}
  Denote the delete operation on the binary min-heap $H$ as $\textsc{Delete-Min}(H)$, which we know from the lectures runs in $O(\log n)$ time, where $n$ is the number of elements in the heap.
  The required algorithm is:
  \begin{algorithm}[H]
    \caption{Check if two binary min-heaps have a common key}
    \begin{algorithmic}[1]
      \State \textbf{Input:} Two binary min-heaps $A$ and $B$ with a total of $n$ elements.
      \State \textbf{Output:} \textbf{True} if $A$ and $B$ have a common key, \textbf{False} otherwise.
      \State Initialize an empty array $a$ of size 2.
      \State $a[0] \gets \textsc{Delete-Min}(A)$
      \State $a[1] \gets \textsc{Delete-Min}(B)$
      \While{$A$ and $B$ are not empty}
        \If{$a[0] = a[1]$}
          \State \Return \textbf{True}
        \ElsIf{$a[0] < a[1]$}
          \State $a[0] \gets \textsc{Delete-Min}(A)$
        \Else
          \State $a[1] \gets \textsc{Delete-Min}(B)$
        \EndIf
      \EndWhile
      \State \Return \textbf{False}
    \end{algorithmic}
  \end{algorithm}

  The algorithm starts by calling \textsc{Delete-Min} on both heaps, taking $O(\log n)$ time (since number of elements in both heaps is less than $n$).
  Each iteration of the while loop calls \textsc{Delete-Min} at most once, again taking $O(\log n)$ time during each iteration.
  And the while loop runs at most $(n - 2)$ times giving us the total time complexity of the algorithm as $O(n \log n) + O(\log n) = O(n \log n)$, while using an array of size 2 as our auxiliary space.

\section{Problem 3}
  Define a new operation on a binary min-heap, \textsc{Delete-Insert}, which deletes and returns the root and inserts a new element at the root.
  The new element is then percolated down to maintain the min-heap property.
  This operation runs in $O(\log n)$ time, where $n$ is the number of elements in the heap.
  The required algorithm is:
  \begin{algorithm}[H]
    \caption{Sort a $k$-sorted array}
    \begin{algorithmic}[1]
      \State \textbf{Input:} A $k$-sorted array $A$ of size $n$.
      \State \textbf{Output:} A sorted array.
      \State Initialize a min-heap $H$ with $A[0], A[1], \dots, A[k-1]$.
      \For{$i = 0 \to (n - k - 1)$}
        \State $x \gets A[i + k]$
        \State $A[i] \gets \textsc{Delete-Insert}(H, x)$
      \EndFor
      \For{$i = (n - k) \to (n - 1)$}
        \State $A[i] \gets \textsc{Delete-Min}(H)$
      \EndFor
      \State \Return $A$
    \end{algorithmic}
  \end{algorithm}

  Where \textsc{Delete-Min} is the standard delete operation on a binary min-heap.
  The algorithm starts by initializing a binary min-heap with the first $k$ elements of the array in an offline manner (since we know the elements of the array), which takes $O(k)$ time (from lectures).
  The first for loop creates a sliding window of size $k$ over the array.
  In the $i$th iteration, we correctly place the $i$th smallest element of the array at index $i$ by calling \textsc{Delete-Insert} with the next element in the array, which takes $O(\log k)$ time.
  This works since each element is displaced by at most $k$ indices and hence will be in the heap during its respective iteration.
  The run time of the first for loop is therefore $O((n - k) \log k)$.
  After there are no more elements left to add to the heap, the second for loop extracts the remaining elements from the heap using the usual \textsc{Delete-Min} operation into the array.
  This will take $O(k \log k)$ time.
  Thus, the total time complexity of the algorithm is $O(k) + O((n - k) \log k) + O(k \log k) = O(n \log k)$.

\section{Problem 4}
  Let $A$ be the given array of size $n$.
  We can easily see that the largest difference between any two elements of the array is $\max(A) - \min(A)$.
  We will use \textsc{Max-Min} (defined later) to find the maximum and minimum elements of an array in $O(n)$ time.
  The required algorithm using this function is simply:
  \begin{algorithm}[H]
    \caption{Find the maximum difference between any two elements of an array}
    \begin{algorithmic}[1]
      \State \textbf{Input:} An array $A$ of size $n$.
      \State \textbf{Output:} The maximum difference between any two elements of $A$.
      \State $(max, min) \gets \textsc{Max-Min}(A, n)$
      \State \Return $max - min$
    \end{algorithmic}
  \end{algorithm}

  We will now define the function \textsc{Max-Min} which takes a divide-and-conquer approach to find the maximum and minimum of an array $A$ of size $n$.
  \begin{algorithmic}[H]
    \Function{Max-Min}{$A, n$}
      \If{$n = 1$}
        \State \Return $(A[0], A[0])$
      \Else
        \State $mid \gets \lfloor n / 2 \rfloor$
        \State $(max_1, min_1) \gets \textsc{Max-Min}(A[0:mid-1], mid)$
        \State $(max_2, min_2) \gets \textsc{Max-Min}(A[mid:n-1], n - mid)$
        \State $max \gets \max(max_1, max_2)$
        \State $min \gets \min(min_1, min_2)$
        \State \Return $(max, min)$
      \EndIf
    \EndFunction
  \end{algorithmic}
  where $A[i:j]$ denotes the subarray of $A$ from index $i$ to $j$ (both inclusive).
  We can easily see that the base case takes $O(1)$ time and each recursive call takes $T(\lfloor n/2 \rfloor) + T(\lceil n/2 \rceil) + O(1) \approx 2 T(n/2) + O(1)$ time (for large $n$).
  To see that the time complexity of the function is $O(n)$ and for simplicity, we can use the substitution $n = 2^k$ and writing $U(k) = T(2^k) = T(n)$ (we can do this because we are analyzing $T$ asymptotically).
  Hence, the recurrence becomes:
  \begin{align*}
    T(2^k) &= 2 T(2^{k-1}) + O(1) \\
    \Rightarrow U(k) &= 2 U(k - 1) + O(1) \\
    &= 2(2 U(k - 2) + O(1)) + O(1) \\
    &= 2^2 U(k - 2) + O(1) + O(1) \\
    &\dots \\
    &= 2^k U(0) + \underbrace{O(1) + \dots + O(1)}_{k \text{ times}} \\
    &= 2^k O(1) + O(k) \\
    &= O(2^k) = O(n) \\
    \Rightarrow T(n) &\in O(n)
  \end{align*}

\end{document}
